\documentclass[10pt,a4paper]{article}

% ----------------- Paquetes -------------------%
\usepackage[T1]{fontenc}
\usepackage[utf8]{inputenc}
\usepackage[a4paper,margin=2.5cm]{geometry}
\usepackage{mathptmx}             
\usepackage{changepage} 
\usepackage{setspace}
\usepackage[spanish]{babel}
\usepackage[labelfont=bf,labelsep=space]{caption}
\usepackage{amssymb}
\usepackage{amsmath}
\usepackage{ebgaramond}
\usepackage{placeins}
\usepackage{etoolbox}
\usepackage{setspace}
\usepackage{parskip} 
\usepackage{tcolorbox}
\usepackage{needspace}
\usepackage{xparse}
\usepackage{ocgx2}
\usepackage{hyperref}
\usepackage{hypcap}


%%%%%%%%%%%%%%%%%%%%%%%%%%%%%%%%%%%%%%%%%%%%%%%%%%%%%%%%%%%%%%%%
% --- Fija primero tus valores globales "buenos" ---
\setstretch{1.2}                 % Interlineado global
\setlength{\parskip}{0.7\baselineskip}
\setlength{\parindent}{0pt}
\newlength{\GlobalParskip}
\setlength{\GlobalParskip}{\parskip}

\newcounter{eqnum}[subsection]
\renewcommand{\theeqnum}{\arabic{eqnum}}
\newcounter{alphaitem}
\newcounter{numitem}

%% ------------------- Recuadros----------------------%%
\newcommand{\puntosclave}[1]{%
  \begin{tcolorbox}[
    colframe=black,
    title={Puntos clave},
    before upper={\setlength{\parindent}{0.5cm}\setlength{\parskip}{0.5\baselineskip}}
  ]
  #1
  \end{tcolorbox}
}

\newcommand{\modificaciones}[1]{
  \begin{tcolorbox}[
    colback=white,
    colframe=red,
    title={Modificaciones a realizar},
    before upper={\setlength{\parindent}{0.5cm}\setlength{\parskip}{0.5\baselineskip}}
  ]
  #1
  \end{tcolorbox}
}

%% ------------------- Listas----------------------%%
    % --- Lista alfabética ---
    \newenvironment{la}
      {\begin{list}{\alph{alphaitem})}{%
          \usecounter{alphaitem}%
          \setlength{\leftmargin}{1cm}%
          \setlength{\parskip}{3pt}% 
          \setlength{\itemsep}{0pt}% ← espacio entre ítems
          \setlength{\parsep}{0pt}%  ← párrafos dentro de un ítem
      }}
      {\end{list}}
      
    % --- Lista numérica ---
    \newenvironment{lnum}
      {\begin{list}{\arabic{numitem}.}{%
          \usecounter{numitem}%
          \setlength{\leftmargin}{1cm}%
          \setlength{\parskip}{3pt}% 
          \setlength{\itemsep}{0pt}% ← espacio entre ítems
          \setlength{\parsep}{0pt}%  ← párrafos dentro de un ítem
      }}
      {\end{list}}
      
% ------------------- Imágenes----------------------%
    \graphicspath{{Img/}}% Rutas por defecto 
    % --- Imagen adaptativa ---
    % Registros de dimensión definidos una sola vez
    \newdimen\imgcap
    \newdimen\imgmin
    \newdimen\imgmax
    \newdimen\imgremain
    \newdimen\imgh
    \newdimen\imgneed
    
    \newcommand{\imagenfit}[3]{%
      \addvspace{10pt}% separación antes
      \par\begingroup
        % Parámetros de composición (asignaciones, no \newdimen)
        \imgcap=1\baselineskip            % alto estimado de título+fuente
        \imgmin=0.15\textheight
        \imgmax=0.15\textheight
        \imgremain=\pagegoal
        \ifdim\pagegoal=\maxdimen \imgremain=0pt\fi
        \advance\imgremain by -\pagetotal
        \advance\imgremain by -\imgcap
        % Decisión de altura destino
        \ifdim\imgremain<\imgmin
          \newpage
          \imgh=0.20\textheight
        \else
          \imgh=\imgremain
          \ifdim\imgh>\imgmax \imgh=\imgmax \fi
        \fi
        % Reservar espacio para evitar cortes del bloque completo
        \imgneed=\imgh \advance\imgneed by \imgcap
        \Needspace{\imgneed}
        % Cuerpo
        \noindent\begin{minipage}{\textwidth}
          \centering
          \captionof{figure}{\textemdash\ #2}\par\vspace{0.3em}% título arriba
          \includegraphics[width=\linewidth,height=\imgh,keepaspectratio]{\detokenize{#1}}\par
          \vspace{0.3em}\small\textit{Fuente: }#3% fuente abajo
        \end{minipage}%
      \endgroup
      \par\addvspace{10pt}%
    }

%------------ Bloque de RECUADRO ESPECIAL -------------%
    \newenvironment{specialblock}[1]{%
      \begin{quote} % crea sangría a izquierda y derecha
      \small % texto más pequeño
      \noindent\textbf{#1}\\[-0.3em] % título en negrita
      \rule{\linewidth}{0.4pt}\\[0.6em]}{
      \end{quote}}
      

%-------------------------- Author Font -----------------------%
    \newcommand{\authorfont}[1]{{\fontfamily{EBGaramond-TLF}\selectfont #1}}

%-------------------------- Anexos -----------------------%
    \makeatletter
    \newcounter{anexo}
    \renewcommand{\theanexo}{\Roman{anexo}}
    \newcommand{\anexo}[1]{%
      \clearpage
      \phantomsection
      \refstepcounter{anexo}%
      \noindent\textbf{Anexo \theanexo.- }#1\par\medskip
      \addcontentsline{stc}{section}{Anexo \theanexo.- #1}%
    }
    \makeatother

    %-------------------------- Lista de figuras (personal) -----------------------%
\makeatletter
\newcommand{\MyListOfFigures}{%
  \clearpage
  \phantomsection
  {\centering\bfseries\Large Índice de figuras\par}%
  \medskip
  % Entrada en nuestro índice de contenido (stc)
  \addcontentsline{stc}{section}{Índice de figuras}%
  % Recuperación del listado (sin cabecera propia)
  \begingroup
    \parindent=0pt\parskip=2pt
    \@starttoc{lof}%
  \endgroup
}
\makeatother

%-------------------------- Bibliografía -----------------------%
\makeatletter
% Contador para las referencias
\newcounter{myref}
% Cabecera de la bibliografía, entra en el índice 'stc'
\newcommand{\MyBibliography}{%
  \clearpage
  \phantomsection
  {\centering\bfseries\Large Bibliografía y referencias\par}%
  \medskip
  \addcontentsline{stc}{section}{Bibliografía y referencias}%
  \setcounter{myref}{0}%
}
\newcommand{\MyBibItem}[4]{%
  \refstepcounter{myref}%
  \noindent[\themyref]\ #1.\ \emph{#2}. #3. #4\par
}
\makeatother

%--------- Establecer cuatro niveles de índice --------%
    \setcounter{tocdepth}{4}   
    \setcounter{secnumdepth}{4}% numera hasta \paragraph

%%%%%%%%%%%%%%%%%%%%%%%%%%%%%%%%

\makeatletter

    %--------- Interlineado Global --------%
    \edef\GlobalBaseLineStretch{\baselinestretch}
    
    %--------- Espaciado de seccinones --------%
    \renewcommand\section{\@startsection{section}
      {1}{\z@}
      {2.5ex \@plus 1ex \@minus .2ex} % espacio antes
      {1.5ex \@plus .2ex}             % espacio después
      {\normalfont\Large\bfseries}    % formato del título
    }
    \renewcommand\subsection{\@startsection{subsection}{2}{\z@}
      {3ex \@plus 0.8ex \@minus .2ex}
      {1.5ex \@plus .2ex}
      {\normalfont\large\bfseries}
    }
    \renewcommand\subsubsection{\@startsection{subsubsection}{3}{\z@}
      {3ex \@plus 0.5ex \@minus .2ex}
      {1ex \@plus .2ex}
      {\normalfont\normalsize\bfseries}
    }
    \renewcommand\paragraph{\@startsection{paragraph}{4}{\z@}%
    {-2ex \@plus -0.5ex \@minus -.2ex}% espacio antes
    {0.8ex \@plus .2ex}% espacio después
    {\normalfont\normalsize\itshape}}% ← cursiva en lugar de negrita

    %--------- Ecuaciones --------%   
    % Variables globales para almacenar títulos y notas
    \gdef\leq@current@section{}%
    \gdef\leq@current@subsection{}%
    \gdef\leq@last@printed@section{}%
    \gdef\leq@last@printed@subsection{}%
    
    % Comando para añadir nota (invisible en el cuerpo)
    \newcommand{\eqnote}[1]{%
      \addcontentsline{leq}{leqnote}{#1}%
    }
    
    % Comando para ecuaciones
    \newcommand{\eqblock}[2]{%
      % Verificar si necesitamos imprimir el título de sección
      \ifx\leq@current@section\leq@last@printed@section\else
        \addcontentsline{leq}{leqsection}{\leq@current@section}%
        \global\let\leq@last@printed@section\leq@current@section
        \gdef\leq@last@printed@subsection{}% Reset subsección al cambiar sección
      \fi
      % Verificar si necesitamos imprimir el título de subsección
      \ifx\leq@current@subsection\leq@last@printed@subsection\else
        \ifx\leq@current@subsection\@empty\else
          \addcontentsline{leq}{leqsubsection}{\leq@current@subsection}%
          \global\let\leq@last@printed@subsection\leq@current@subsection
        \fi
      \fi
      % Ecuación
      \refstepcounter{eqnum}%
      \begin{equation*}
        #1\tag{E.\theeqnum}
      \end{equation*}%
      \addcontentsline{leq}{leqt}{[E.\theeqnum] -- #2}%
      \addcontentsline{leq}{leqm}{#1}%
    }
    
    %%% ----- Índice de ecuaciones ----------%%%
    % Tamaño para []
    \newcommand{\leqIndexFont}{\normalsize}
    \newcommand{\leqTitleFont}{\tiny}
    \newcommand{\leqNoteFont}{\tiny}
    % Cabecera de sección 
    \newcommand*\l@leqsection[2]{%
      \addvspace{25pt}% Mayor separación antes de nueva sección
      \noindent\leqIndexFont
      \makebox[\linewidth]{\rule{.38\linewidth}{0.4pt}\ \ \textbf{  \MakeUppercase{#1}}\ \ \rule{.38\linewidth}{0.4pt}}%
      \par\vspace{-10pt}
    }
    % Cabecera de subsección 
    \newcommand*\l@leqsubsection[2]{%
      \par\addvspace{15pt}%
      \noindent\leqIndexFont #1
      \par\vspace{-7pt}
    }
    % Título de ecuación 
    \newcommand*\l@leqt[2]{%
    \par\addvspace{10pt}%
      \par\noindent\leqTitleFont #1\par
    }
    % Miniatura de ecuación
    \newcommand*\l@leqm[2]{%
      \vspace{0.3\baselineskip}%
      \par\scriptsize\noindent\hspace*{1.5em}\(#1\)\par%
    }
    % Nota de ecuación 
    \newcommand*\l@leqnote[2]{%
      \vspace{0.2\baselineskip}%
      \par\noindent\leqNoteFont\hspace*{1.5em}\textbf{Nota.--} #1\par\vspace{0.5\baselineskip}%
    }
    % Generación del índice
    \newcommand{\mylistofequations}{%
      \clearpage
      \phantomsection
      % Título visual de la lista (ajústalo a tu gusto):
      {\centering\bfseries\Large Índice de ecuaciones\par}
      % Entrada en nuestro índice 'stc':
      \addcontentsline{stc}{section}{Índice de ecuaciones}%
      % Llamada a tu lista real (usa el nombre exacto de tu macro):
      \@starttoc{leq}%
    }
    
    %--------- Captura de títulos de sección --------%
    \let\leq@original@section\section
    \let\leq@original@subsection\subsection
    % Flag para saber si estamos en una sección asterisco
    \newif\ifLeqStarSection
    \LeqStarSectionfalse
    
    % Redefinir \section CON CONTROL DE ASTERISCO
    \renewcommand{\section}{%
      \@ifstar{\leq@section@star}{\@dblarg\leq@section@nostar}%
    }
    \def\leq@section@star#1{%
      \global\LeqStarSectiontrue
      \gdef\leq@current@section{#1}%
      \gdef\leq@current@subsection{}%
      \leq@original@section*{#1}%
      \global\LeqStarSectionfalse
    }
    \def\leq@section@nostar[#1]#2{%
      \global\LeqStarSectionfalse
      \gdef\leq@current@section{#2}%
      \gdef\leq@current@subsection{}%
      \leq@original@section[#1]{#2}%
    }
    % Redefinir \subsection
    \renewcommand{\subsection}{%
      \@ifstar{\leq@subsection@star}{\@dblarg\leq@subsection@nostar}%
    }
    \def\leq@subsection@star#1{%
      \global\LeqStarSectiontrue
      \gdef\leq@current@subsection{#1}%
      \leq@original@subsection*{#1}%
      \global\LeqStarSectionfalse
    }
    \def\leq@subsection@nostar[#1]#2{%
      \global\LeqStarSectionfalse
      \gdef\leq@current@subsection{#2}%
      \leq@original@subsection[#1]{#2}%
    }

    % ----- Restauración de valores Globales de estilo ----------%
    % Flags
    \newif\ifInFootnote
    \InFootnotefalse
    \pretocmd{\@footnotetext}{\InFootnotetrue}{}{}
    \apptocmd{\@footnotetext}{\InFootnotefalse}{}{}
    % Profundidad de listas (soporta anidadas)
    \newcount\MyListDepth \MyListDepth=0
    \AtBeginEnvironment{la}{\advance\MyListDepth by 1}
    \AfterEndEnvironment{la}{\advance\MyListDepth by -1}
    \AtBeginEnvironment{lnum}{\advance\MyListDepth by 1}
    \AfterEndEnvironment{lnum}{\advance\MyListDepth by -1}
    \AtBeginEnvironment{itemize}{\advance\MyListDepth by 1}
    \AfterEndEnvironment{itemize}{\advance\MyListDepth by -1}
    \AtBeginEnvironment{enumerate}{\advance\MyListDepth by 1}
    \AfterEndEnvironment{enumerate}{\advance\MyListDepth by -1}
    % Restauración diferida: hazlo al INICIO del próximo párrafo real
    \newif\if@pendingRestore
    \@pendingRestorefalse
    \newcommand{\RequestRestore}{\global\@pendingRestoretrue}
    \everypar{%
      \if@pendingRestore
        \global\@pendingRestorefalse
        \ifInFootnote\else
          \ifnum\MyListDepth=0
            \setlength{\parskip}{\GlobalParskip}%
            \setstretch{\GlobalBaseLineStretch}%
          \fi
        \fi
      \fi
    }
    % Al final de bloques:
    \AfterEndEnvironment{equation}{\RequestRestore}
    \AfterEndEnvironment{equation*}{\RequestRestore}
    \AfterEndEnvironment{la}{\RequestRestore}
    \AfterEndEnvironment{lnum}{\RequestRestore}
    % Al final de macros:
    \apptocmd{\eqblock}{\RequestRestore}{}{}
    \apptocmd{\imagenfit}{\RequestRestore}{}{}
    
\makeatother

% ===================== ToC propio con 4 niveles (único bloque) =====================
\makeatletter

% --- Escritores: añadimos una línea al .stc en cada cabecera numerada ---
% Guardamos las versiones actuales para llamar al comportamiento normal
\let\MyTOC@orig@section\section
\let\MyTOC@orig@subsection\subsection
\let\MyTOC@orig@subsubsection\subsubsection
\let\MyTOC@orig@paragraph\paragraph

% SECTION
\renewcommand{\section}{%
  \@ifstar{\MyTOC@section@star}{\@dblarg\MyTOC@section@nostar}%
}
\def\MyTOC@section@star#1{%
  \MyTOC@orig@section*{#1}% sin entrada al .stc
}
\def\MyTOC@section@nostar[#1]#2{%
  \MyTOC@orig@section[{#1}]{#2}% comportamiento normal (escribe al .toc)
  % Escribimos también nuestra línea al .stc (texto con \numberline)
  \addcontentsline{stc}{section}{\protect\numberline{\thesection}#2}%
}

% SUBSECTION
\renewcommand{\subsection}{%
  \@ifstar{\MyTOC@subsection@star}{\@dblarg\MyTOC@subsection@nostar}%
}
\def\MyTOC@subsection@star#1{%
  \MyTOC@orig@subsection*{#1}%
}
\def\MyTOC@subsection@nostar[#1]#2{%
  \MyTOC@orig@subsection[{#1}]{#2}%
  \addcontentsline{stc}{subsection}{\protect\numberline{\thesubsection}#2}%
}

% SUBSUBSECTION
\renewcommand{\subsubsection}{%
  \@ifstar{\MyTOC@subsubsection@star}{\@dblarg\MyTOC@subsubsection@nostar}%
}
\def\MyTOC@subsubsection@star#1{%
  \MyTOC@orig@subsubsection*{#1}%
}
\def\MyTOC@subsubsection@nostar[#1]#2{%
  \MyTOC@orig@subsubsection[{#1}]{#2}%
  \addcontentsline{stc}{subsubsection}{\protect\numberline{\thesubsubsection}#2}%
}

% PARAGRAPH
\renewcommand{\paragraph}{%
  \@ifstar{\MyTOC@paragraph@star}{\@dblarg\MyTOC@paragraph@nostar}%
}
\def\MyTOC@paragraph@star#1{%
  \MyTOC@orig@paragraph*{#1}%
}
\def\MyTOC@paragraph@nostar[#1]#2{%
  \MyTOC@orig@paragraph[{#1}]{#2}%
  % Si no numeras los \paragraph, cambia \theparagraph por texto plano sin \numberline
  \addcontentsline{stc}{paragraph}{\protect\numberline{\theparagraph}#2}%
}

% --- Impresor del índice propio (lee el .stc) ---
\newcommand{\MyTOC}{%
  \par\bigskip
  {\centering\bfseries\Large Índice de contenido\par}%
  \medskip
  \begingroup
    \parindent=0pt\parskip=2pt
    \@starttoc{stc}%
  \endgroup
}

\makeatother
% ===================== fin ToC propio =====================


%%%%%%%%%%%%%%


% --- Datos del tema ---
\title{Teoría impositiva (I): Ingresos públicos\\
\large Elementos definidores y clases de impuestos. Principios impositivos}
\author{Marcelo Ortega}
\date{Agosto 2026}

\usepackage{soul}\sethlcolor{yellow}% resaltado amarillo: \hl{texto} (Panel TCEE, ⌃H)
\begin{document}
\maketitle
\newpage

% --- Página de resumen y modificaciones ---
\puntosclave{ %% Escribir puntos clave Apuntes CECO
     Este es un tema corto debido a que el contenido da para pocas ampliaciones sin invadir el contenido de otros temas. La cuestión clave es ser sistemático y tratar de otorgar el tiempo necesario a cada apartado del mismo. En este sentido, se puede dividir el tema en cuatro partes que pueden ser tratadas de igual manera (todos de \textbf{siete minutos}):
     \begin{lnum}
         \item \textbf{Definición de ingresos públicos}. Parte ampliada sobre los ingresos públicos de acuerdo al SEC 2010 y al PGE que sigue la contabilidad pública. Esta primera parte debe ser concisa pero no debe recortarse. Existen diferentes tipos de ingresos y aquí es donde deben dejarse claros;
         \item \textbf{Elementos definidores: impuestos}. Los impuestos como principal fuente de ingresos de las economías desarrolladas tienen una serie de elementos definidores. Estos elementos definidores son comunes a todos los tributos y constituyen una liquidación. Por ello, se recomienda llevar a cabo una \textit{liquidación tipo} esquemática que facilite la exposición oral;
         \item \textbf{Clases de impuestos}. Las clases de impuestos se pueden definir según los criterios anteriores. Esta parte es la que corre más riesgo de quedarse corta y se anima a ampliarla con conocimientos específicos de cada tributo una vez se entre en la clasificación directo/indirecto. Ejemplificar cada tipo de tributo de acuerdo a esta clasificación. Lo más importante es contar con los conocimientos de otros temas para cumplir con el tiempo;
         \item \textbf{Principios impositivos}. Esta parte de economía política y debate se pone al final pues puede ayudar a rellenar falta de tiempo. La clave de los principios impositivos es que estos se deben analizar para el sistema tributario en su conjunto y no sólo para un tributo.
     \end{lnum}
    }
\vspace{1cm}

\modificaciones{
    
    }

\newpage
\MyTOC
\newpage

\mylistofequations
\clearpage

\MyListOfFigures
\clearpage


\pagenumbering{arabic}
\setcounter{page}{1}


\section{Introducción}
En este tema analizaremos los ingresos públicos como fuente de recursos para las Administraciones Públicas (sea cual sea el subsector) y que permite llevar a cabo las funciones de la economía pública ele ordenación y planificación, facilitando la corrección de errores de mercado y mejorando la redistribución y la equídad, en un contexto de una economía mixta. Es por ello que el concepto de Sistema de ingresos debe quedar remarcado, pues los ingresos de las Administraciones Públicas no deben analizarse de manera aislada. es decir tributo a tributo o fuente a fuente de ingresos. si no como un conjunto o sistema impositivo o de ingresos. De esta manera se es capaz de analizar las relaciones de complementariedad y sustituibilidad que existen entre las distintas fuentes de ingresos y a su vez los impactos que el hecho de elegir una fuente uotra implican.


\subsection{Contextualización}


\subsubsection{Relevancia}




\subsubsection{Problemática}



\subsection{Estructura de la exposición}





\clearpage
\section{Ingresos públicos: Marco conceptual}
Los ingresos públicos se pueden definir como toda cantidad de recursos obtenidos por las Administraciones Públicas, con el doble objetivo de:
\begin{la}
    \item \textbf{Financiar}. La financiación del gasto público, como principal herramienta para la consecución de los objetivos y fines marcados por la Administración Pública; y,
    \item \textbf{Distorsionar}. Modificar el comportamiento de los agentes económicos asociada al proceso de recaudación del ingreso, como herramienta para la consecución de objetivos marcados por la Administración Pública (fumar).
\end{la}

La naturaleza de los ingresos públicos es muy diversa dependiendo de la estructura económica de cada país. Algunos países obtienen sus ingresos mayoritariamente de impuestos (España) otros de recursos naturales (Noruega, Guinea Ecuatorial, etc.), otros de la buena voluntad de sus vecinos que le transfieren recursos (Somalia, Sudan del Sur).

En definitiva, los ingresos públicos son muy diversos y este tema debe centrarse en analizar estos desde una óptica universal y no sólo aquella que puede afectar a un país desarrollado con economía mixta. Es por ello que para analizar los tipos de ingresos recurriremos a la clasificación de Naciones Unidas.

\subsection{Clasificación de los ingresos: Sistema de Cuentas Nacionales (SCN, 2008)}
Naciones Unidas elabora el Sistema de Cuentas Nacionales 2008 (SCN 2008) que es un marco estadístico que proporciona un conjunto completo, coherente y flexible de cuentas macroeconómicas para la formulación de políticas, análisis y propósitos de investigación. Se ha producido y está realizado bajo los auspicios de las Naciones Unidas. la Comisión Europea, la Organización para la Cooperación y el Desarrollo Económico. el Fondo Monetario Internacional y el Banco Mundial. La Unión Europea tiene su correlato equivalente que es el Sistema Europeo de Cuentas 2010 (SEC 2010). Este sistema clasifica los ingresos públicos de manera económica. lejos de las identidades nacionales y de manera universal para todas las economías del planeta.

Así la clasificación posible de ingresos es la siguiente:
• Impuestos sobre la producción y las impo1taciones
o impuestos sobre l a producción y las importaciones
o impuestos sobre los productos
• Impuestos corrient(ls sobre el ingreso, la riqueza etc.
o Impuestos sobre los hogares y otros.
o Impuestos sobre las empresas
• impuestos sobre el capital
• Contribuciones sociales netas
• Renta de la propiedad, ventas y otros ingresos.
• Otras transferencias de capital

\subsubsection{Impuestos sobre la producción y las importaciones}
Son impuestos sobre la actividad independientemente de que esta genera renta.
Su definición oficial4 es:
\textit{'·Los impuestos sobre la producción y las importaciones (D.2) son pagos obligatorios sin contrapartida, en efectivo o en especie, recaudados por las administraciones públicas u por lus ínstifuciones de la Unión Europea. que gravan fa producción y la importación de bienes y servicios, la 11/ilizac:íón de mano de obra. la propiedad u el uso de la fierra, los edificios v 01rosacrivos utili:::ados en la producció11. Dichos ímpuesros deben pagarse. i11dependientemenfe de los beneficios ohrenidos. ,.}
Se dividen en:
• impuestos sobre los productos
o impuestos del tipo de valor añadido (IV A)
o impuestos y derechos sobre las importaciones, excluido el JV A
o impuestos sobre los productos, excluidos el IV A y los impuestos sobre las importaciones
• otros impuestos sobre la producción.

Los impuestos sobre los productos Son impuestos a pagar por cada unidad de un dete1minado bien o servicio producido o negociado. El impuesto puede consistir en un montante monetario específico por unidad de un bien o servicio, o puede calcularse como un porcentaje específico del precio unitario o del valor de los bienes y servicios producidos o negociados. Los impuestos que gravan un producto, independientemente de la unidad institucional que los pague, se incluyen en los impuestos sobre los productos, a menos que estén incluidos específicamente en otra rúbrica (IV A, llEE).

Los impuestos sobre la producción comprenden todos los impuestos que soportan las empresas como resultado de su pat1icipación en la producción, independientemente de la cantidad o el valor de los bienes y servicios producidos o vendidos (1131, IAE).

\subsubsection{Impuestos corrientes sobre el ingreso y la riqueza}
Son impuestos que gravan mani festaciones de riqueza asociadas a rentas o riqueza. Es por ello que su naturaleza es próxima a la ele concepto de impuestos directos.
\textit{·'Lo.1 impuestos corri1?1 1tes sohrl! la rent a. el patr imonio. etc . . comprenden todo.1 los pagos obliga1urios sin contrapurlida. en t!/t !ctivu o <!11 t!Spe<.:ie. n!caudwlos ¡1eriúdicwne111e por las administracio nes públicas y por el resto del mundo sohrl! la renta v t!l palrimonio di! las unidadt!s in sliluciunales. así comu alguno.- impuesto.\ periódicos que nu se exigl!n ni sobre dicha renta ni sobre dicho ¡1a1rimo11io . • ,}

Los impuestos corrientes sobre la renta. el pat1imonio, etc., se dividen en:
• Impuestos sobre la renta, y estos a su vez en:
o impuesto sobre las personas e IFHSL e
o impuesto sobre las sociedades
• Otros impuestos co1Tíentes

Los primeros son los de más peso. y en ellos enmarcaríamos las fi guras nacionales de Impuesto sobre la Renta de las Personas Físicas e Impuesto sobre Sociedades. El ��egundo es una categoría residual y menor. donde encuad raríamos tributos nacionales como el Impuesto sobre el Patrimonio o Impuesto sobre Vehículos de Tracción Mecánica

\subsubsection{Impuestos sobre el capital}
Son impuestos que gravan manifestaciones puntuales de riqueza, asociadas generalmente a una transferen cia de patrimonio o riqueza. La clave de estos impuestos es que recaen sobre manifestaciones recurrentes ele riqueza y no sobre rentas.

'·Los impuestos sobre el capit al son impuestos que f{rava, 1 a i111ervalos irregulares y 11111v poco ji-ecuentes el valor de los activos o el patrimonio neto de las unidade s institucionales o e l valor de los activos transferidos entre unidades institucion ales com o resultado de sucesiones, donaciones inter vivos u o/ras tromferenci as. " 

Dentro de esta figura encuadraríamos los Impuestos de Sucesiones y· Donaciones y el Impuesto sobre el Incremento del valor de los TeJTenos de Naturaleza Urbana (Plusvalía).

\subsubsection{Contribuciones sociales netas}
Recogen todas las contribuciones sociales pagadas dentro del sistema de Seguridad Social, sea por empleador o empleado, sean efectivas o imputadas.
'·Las cotizaciones sociales netas son las coti:::aciones efectivas o imputadas que pagan los hogares a los sisiemas de seguros sociales con el.fin de asegurar el pago de pres1acione.1· s ociales.''

Estas a su vez son finalistas y sirven para cubrir contingencias a las que se enfrentan los hogares, como pueden ser el desempleo o la enfem1 edad.

'·L as prestaciones sociales s on tran.��ferencias a los hogares, en efectivo o en esp ecie, de stinadas a aliviar la carga.fi11a11ciera que representa para ellos la cobert ura de una serie de riesgos o necesidades y ef ectuadas por medio de sist1?111 as mxanizados colectivamente o, fuera de estos sistemas . por 1111idades de las administraciones públicas y por las ISFLSH. Dichas transferencias incluyen los pagos de las administraciones públicas a los productores que suministran prestaciones de forma individual a los hogares en el marco de los riesgos y necesidades sociales. ··

En este caso se suele diferenciar entre imputadas y efectivas, la diferencia esencialmente es que las imputadas se asocian a prestaciones que se dan directamente por parte del empleador sin mediar sistema de protección social (por ejemplo, ayudas de acción social o ayudas a dentistas que detemlinadas empresas pagan). Las contribuciones sociales imputadas serán más importantes en sistemas como el de los Estados Unidos.

\subsubsection{Renta de la propiedad, ventas y otros ingresos}
·'Las rentas de la propiedad se generan cuando los propietarios de activos financieros y de recursos 1,aturales los ponen a di.sposición de otrus unidades institucionales. La rento a pagar por la utilización de activosj,11a11cieros se denomina rentas de inversión. mientras que la que se paga por la utilización de w1 recurso na111ral se denomina alquiler. Las rentas de la propiedad son la suma de las rentas de inversión más los alquileres. "

Esta categoría de ingresos, si bien no de mucha importancia en países desarrollados (y con muchas salvedades, véase Noruega) supone una fuente de renta importante para muchos países y va asociada a su riqueza i ntrínseca o geográfica de manera directa, como una concesión o alquiler o indirecta, como un fondo soberano. Podría llegar a darse el caso de una nación que a base de acumular superávit constituyese un fondo soberano y obtuviese rentas de él; de momento el autor no tiene conocimiento de que dicho caso se haya dado, de ahí su relación con la riqueza intrínseca del país.

'·Las rentas de la propiedad se generan cuando los propietarios de activos financieros )' de recursos naturales los ponen a disposición de otras 1111idades inslitucionales. La renta a pagar por la utilización de activosfinancieros se denomina rentas de inversión, mientras que la que se paga por la utili=ación de un recu rso natural se denomina alquiler. Las rentas de la propiedad son la suma de las l"l?ntas de inversión más los alquileres :·

Otras transferencias de capital.

"'Las otras transferencias de capital comprenden las tran4erencias distintas de las ayudas a la inversión y los impuestos sobre el capital, q11e no redistribuyen por sí mismas la renta, pero que sí comportan una redistribución del ahorro o del patrimonio emre los distintos sectores u subsectores de la economía o entre estos y el resto del mundo. Pueden efectuarse en efectivo o en <?specie (casos de as11nción de deuda o de cancelación de deuda) y curresponden a tran��fere11cias volunlarias de riqU<?za. ,.

Es una fuente de in¡,,-resos unilateral que puede provenir de otro subsector o del exterior sin contraprestación. Un ejemplo es determinados tipos de transferencias de la UE o ele alguno organismo internacional.

El conjunto de fuentes de ingreso sería lo que definiría los ingresos públicos en un sentido amplio de acuerdo a los Sistemas de Cuentas Nacionales de Naciones Unidas. Estos son aplicables a cualquier país y no dependen de la estructura impositiva propia del país. Indicadores de ingresos tributa1ios.

La UE dentro del marco que establece el SEC 2 0 1 O rea liza una serie de definiciones alternativa y de conjuJJtos de indicadores que pasamos analizar. De acuerdo a Eurostat los ingresos tributarios son el total de impuestos y contribuciones sociales que se pagan al gobierno general (conjunto de todos los subsectores), incluidas aquellas aportaciones voluntarias. El hecho de contar con un ind icador es de gran importancia pues servirá para calcular métricas que ayuden a la toma de decisiones de orientación de política fiscal:

Presión fiscal=lngresos tributarios/Pll3,

Esta definición la más amplia será la que acoja el cuarto indicador que usa la UE para definir los ingresos.\footnote{%
    Entiéndase esta clasificación como próxima conceptualmente a la definición de tipos de dinero.
}
.
+ Im puestos sobre la producción y las impo11aciones
+ Impuestos conientes sobre el ingreso, la riqueza etc.
+ Impuestos sobre el capital
- transferencias de capital asociadas a impuestos y contribuciones sociales de dudoso cobre
+ contribuciones sociales obligatorias pagables a fondos de seguridad social .
= Indicador 1
+ contribuciones sociales obligatorias pagables a todos los subsectores del gobierno menos fondos
de seguridad social
= Indicador 2
+ contribuciones sociales imputadas al gobierno como empleador
+ complementos de los hogares a las contribuciones sociales
- costes de servicio de los sistemas de seguro
= Indicador 3
+ contribuciones sociales vol untarias.
= Indicador 4

Esta c lasificación de Eurostat acomoda todos los posibles sistemas tributarios de la UE, con especial atención a las contribuciones sociales. que presentan importantes peculiaridades. Sirva corno ejem plo para i lustrar las fuentes de ingresos en la UE el siguiente gráfico:

\subsection{Clasificación: Contabilidad Pública}
NOTA: Antes de continuar la parte restante del tema, que se centra exclusiwnnente en Impuestos. creemos necesario hacer una mención u como se clasifican los ingresos públicos en E.1pa1ia. de acuerdo a la contabilidad pública. De i?sta forma esta clasificacián será la que se enrnentre en los f'resupuesro:; Ge11erales del Gobiemo y la anterior, la que encontraremos en documentos como el Programa de Estabilidad, Comunicaciones sobre Déficit Excesivo y demás documentosque podríamos denominar "externos ".

Mantenemos nuestra definición de ingreso público como "Toda cantidad de recursos percibida por el Gobierno cuyo objetivo principal es· la financiación del gasto público, como principal herramienta de los objetivos y fines del m i smo."

Y los clasificamos en dos categorías:

• Ingresos tributarios. Estos serán las cantidades que el Sector Público detrae coactivamente del Sector Privado ( Contribuyentes), generalmente bajo mandato legislativo. En democracias se suele sustentar en el mandato constitucional bajo el principio ele apo,tación al bien común. Cabe que los ingresos tributarios sean de dos tipos:
o Con contraparti da: Estos serían en los que se recibe alguna prestación: tasas o precios públicos.
o Sin contrapartida: no se adquiere derecho a mayores prestaciones por el grado de contribución: Impuestos, cotizaciones sociales.

• Ingresos no tributarios. Estos se obtienen de contribuyentes del sector privado los cuales trasfieren una renta en virtud de una actividad desarrollada por el sector público, ya sea para obtener fondos o en el desarrollo de sus funciones. Esencialmente se trata de ingresos procedentes de rentas de la propiedad o patrimonio del Estado, organismos autónomos, agencias estatales y otros organismos públicos, así como los derivados de activi dades realizadas en régimen de derecho privado. Las más relevantes pueden ser las primas de reembolso de deuda pública o los ingresos del Banco de Esparía por Señoreaje.

Sin embargo, en Espafia la norma, la Ley G?neral Tributaria. indica que existen tres tipos de tributos. Esta clasificación si bien no es universal. como las anteriores, es de utilidad para el opositor para encmiclrar los posteriores temas de tributos espafioles.

• Tasas son tributos cuyo hecho imponible consiste en la obtención de un servicio o bien por parte del sujeto pasivo. Es decir que el hecho imponible es la utili;,ación rrivativa o el aprovechamiento del dominio público: así como la prestación de servicio�� o la realinción de actividades en régimen de derecho público que beneticien de modo particular al obligado tributario o al conj unto de obligados tributarios. Esto será así si los servicios o actividades no son de solicitud o recepción voluntaria para los obligados tributarios o no es posible que se presten o realicen por parte del sector privado. La clave en su detinición es que es el su_ie10 pasi\'O el que la promueve. pero sólo l a administración pública puede prestarlo.

• Contribuciones especiales son los tributos cuyo hecho imponible consiste en la obtención por el obligado 1ributa1io de un benelicio o de un aumemo ele valor de sus bienes como consecuencia de la realización de obras públicas o del establecimiento o ampliación de servicios públicos. En la contribución especial existe también una actividad administrativa. pero, a diferencia de la tasa, dicha actividad surge sin que medie una petición del  contribuyente. es decir. existe una actuación de la Administración que. sin ir dirigida a un sujeto pasivo en especial, le ocasiona un beneficio particular e identificable y por ello debe pagar el tributo. Por e_iemplo, si a un ba1Tio se le dota de un acceso por la autopista, se podría cargar una contribución a los residentes de ese barrio.

• Impuestos son los tributos exigidos sin contraprestación. he aquí la principal diferencia con las dos figuras anteriores. El hecho imponible está constituido por una manifestación de la capacidad económica del contribuyente. Frente a los tipos anteriores. el elemento objetivo del hecho imponible no supone una actividad administrativa y la realización del hecho imponible corresponde excl usiva y únicamente al sujeto pasivo. sin intervención alguna de la Administración, apareciendo la coercitividad en el momento del nacimiento de la obligación tributatia.





\clearpage
\section{Impuestos: Marco conceptual}



Al asumir el cargo en enero de 2001, el presidente George W. Bush se encontró ante una oportunidad poco común: un superávit presupuestario federal proyectado. Bush decidió rápidamente que el gobierno debería utilizar este superávit para reducir la carga fiscal federal. Como dijo en un discurso radiofónico el 24 de febrero de 2011: «Un superávit en los ingresos fiscales, después de todo, significa que a los contribuyentes se les ha cobrado de más. Y normalmente, cuando se te ha cobrado de más, esperas que te devuelvan algo».

Más tarde, durante su primer año, Bush promulgó la Ley de Reconciliación de Crecimiento Económico y Alivio Fiscal de 2001. Esta ley redujo los tipos impositivos de manera generalizada, pasando los tipos del 15 al 39,6 % al 10 al 35 %. El impuesto sobre sucesiones aplicado a los patrimonios de los fallecidos más ricos de Estados Unidos fue eliminado gradualmente, antes de desaparecer por completo en 2010. Estas reducciones fueron ampliadas aún más por la Ley de Reconciliación de Alivio Fiscal para el Empleo y el Crecimiento de 2003, que aceleró el ritmo al que se introducían gradualmente los tipos impositivos más bajos, redujo el tipo impositivo máximo sobre las rentas procedentes de dividendos de acciones y ganancias de capital, y aumentó las deducciones del impuesto de sociedades.

Desafortunadamente, como se analizó en el Capítulo 4, para cuando el presidente Bush dejó el cargo, este superávit se había convertido en un déficit muy grande. Un debate continuo durante toda la posterior presidencia de Obama fue si poner fin a las reducciones fiscales de Bush y cómo hacerlo. Estaba previsto que las reducciones expiraran en 2010, y Obama propuso ampliar las reducciones durante al menos otro año para quienes ganaran menos de 250.000 dólares, lo que incluye al 98 % de los estadounidenses, pero volver a elevar los impuestos a sus niveles originales para quienes ganaran más.

Obama explicó: «Estas reducciones fiscales para los estadounidenses más ricos son también las reducciones fiscales que tienen menos probabilidades de promover el crecimiento».² Pero los republicanos argumentaron que unos impuestos más altos en medio de una recesión perjudicarían aún más a la economía al castigar a las pequeñas empresas creadoras de empleo que pagan impuestos conforme al código del impuesto sobre la renta individual.

Al final, las dos partes llegaron a un compromiso ampliando el conjunto completo de reducciones fiscales durante otros dos años, junto con otras desgravaciones fiscales y gasto para promover el crecimiento durante la recesión en curso. El presidente Obama prometió no permitir otra ampliación de este tipo, diciendo: «He dicho antes que sentía que las reducciones fiscales para la clase media estaban siendo tomadas como rehenes por las reducciones fiscales para las rentas altas. Creo que resulta tentador no negociar con quienes toman rehenes, a menos que el rehén resulte perjudicado. Entonces la gente cuestionará la sensatez de esa estrategia. En este caso, el rehén era el pueblo estadounidense, y yo no estaba dispuesto a permitir que resultara perjudicado».³ No estaba claro si la promesa de no permitir otra ampliación de las reducciones fiscales de Bush era creíble, dado el control republicano del Congreso.

Cuando las reducciones fiscales expiraron (de nuevo) en 2012, se alcanzó otro compromiso mediante la Ley de Alivio para los Contribuyentes Estadounidenses de 2012. Esta ley hizo permanentes todas las reducciones fiscales para aquellos hogares que ganaban menos de 450.000 dólares; por encima de ese nivel, las reducciones fiscales expiraron y el tipo impositivo máximo aumentó del 35 al 39,6 %. Las deducciones fiscales y diversos créditos fiscales también comenzaron a eliminarse gradualmente para quienes ganaban más de 250.000 o 300.000 dólares, para individuos y parejas, respectivamente. El impuesto sobre sucesiones también aumentó del 35 al 40 %, pero siguió manteniendo la misma exención de 5 millones de dólares (es decir, solo se pagaba el impuesto sobre sucesiones por cualquier cantidad superior a 5 millones de dólares).

Estos debates ponen de relieve el importante papel que desempeñan los impuestos tanto en el ámbito político como en la formulación de políticas gubernamentales en Estados Unidos. En esta sección del libro, vamos más allá de nuestro estudio de los gastos públicos para estudiar los ingresos públicos recaudados mediante la imposición. Este capítulo comienza el estudio de la imposición estableciendo el marco institucional y teórico para comprender la política fiscal y sus efectos. Una vez que comprendamos estos conceptos básicos, nuestro estudio de la imposición procede en tres pasos. Los Capítulos 19 y 20 abarcan la teoría básica de la imposición. Los Capítulos 21 a 24 aplican esta teoría básica al estudio de cómo los sistemas fiscales afectan al comportamiento económico de los individuos y de las empresas. Finalmente, el Capítulo 25 analiza las implicaciones de nuestro análisis para una reforma fiscal fundamental en Estados Unidos.

Comenzamos nuestro estudio proporcionando una breve visión general de los tipos de imposición que existen en Estados Unidos, en los niveles federal, estatal y local, y en todo el mundo. A continuación, analizamos con más detalle la principal herramienta de recaudación de ingresos en Estados Unidos: el impuesto federal sobre la renta. Proporcionamos una visión general de la estructura del impuesto sobre la renta y analizamos medios alternativos de medir la «equidad» del sistema del impuesto sobre la renta. A continuación, pasamos a la cuestión de cómo medir la base sobre la que deberían aplicarse los impuestos sobre la renta. ¿Deberían gravarse todas las formas de renta o solo algunas formas? ¿Debería el gobierno utilizar el código fiscal para fomentar la provisión privada de bienes públicos eximiendo de impuestos las contribuciones destinadas a dichos bienes? ¿Debería el gobierno gravar a los individuos sobre la base de su propia renta o sobre la base de la renta de todos los miembros de su familia? Esta parte del capítulo analiza cómo deberían los gobiernos, y cómo lo hacen en la práctica, formar las bases para gravar la renta.


Tipos de imposición

Los gobiernos de Estados Unidos y de otras naciones recaudan ingresos mediante una amplia variedad de mecanismos. Nuestro estudio de la imposición se centrará en los cinco tipos de impuestos más comunes, que revisamos en esta sección. (Otros tipos de impuestos más específicos también se analizarán, cuando corresponda, en los capítulos restantes del libro).

Impuestos sobre los ingresos del trabajo

El primer tipo de imposición es el impuesto sobre las nóminas, un impuesto aplicado a los ingresos de los trabajadores. Los impuestos sobre las nóminas son el principal medio de financiación de los programas de seguro social, como los analizados en los capítulos anteriores (Seguridad Social, seguro de desempleo, Medicare, etc.).

Impuestos sobre la renta individual

El segundo tipo de imposición es el impuesto sobre la renta individual, un impuesto pagado por los individuos sobre la renta acumulada durante el año. La renta a efectos del impuesto sobre la renta incluye los ingresos del trabajo, pero el impuesto se distingue del impuesto sobre las nóminas porque (1) se aplica a un conjunto más amplio de fuentes de renta (como también los ingresos por intereses procedentes del ahorro de los hogares) y (2) se aplica en muchos casos a la renta total de una familia, no solo a la renta de un trabajador individual. Una forma de imposición sobre la renta que reviste especial interés es la imposición de las ganancias de capital, los ingresos procedentes de la venta de activos de capital, como acciones, pinturas y casas.

Impuestos sobre la renta de las sociedades

Además de gravar la renta individual, muchos países también gravan por separado los ingresos de las sociedades mediante el impuesto sobre la renta de las sociedades. El propósito de la imposición separada de las sociedades, por encima y más allá de los impuestos sobre los individuos, es gravar los ingresos de los propietarios del capital que, de otro modo, podrían escapar a la imposición mediante el sistema del impuesto sobre la renta basado en los individuos.

Impuestos sobre la riqueza

Los impuestos sobre la riqueza son impuestos pagados no sobre la renta a medida que se acumula, sino sobre el valor de los activos poseídos por una persona o familia, como terrenos, joyas, obras de arte, bienes inmuebles y acciones. Se incluyen en esta categoría los impuestos estatales y locales sobre la propiedad, que se basan en el valor del terreno y de cualquier estructura construida sobre él, y los impuestos sobre sucesiones, que se basan en los legados (dinero, propiedades, etc.) que se dejan cuando uno muere.

Impuestos sobre el consumo

La forma de imposición más común en todo el mundo es el impuesto sobre el consumo, que se paga sobre el consumo individual o familiar de bienes (y a veces también sobre servicios). Los impuestos sobre el consumo se aplican a menudo en forma de impuestos sobre las ventas, impuestos que los consumidores pagan a los vendedores en el punto de venta. Estos impuestos pueden aplicarse bien a una amplia variedad de bienes de consumo o bien únicamente a un bien concreto. Cuando se aplican solo a determinados bienes, por ejemplo cigarrillos, automóviles o gasolina, el impuesto sobre las ventas se denomina impuesto especial.

Los impuestos sobre las nóminas, sobre la renta y sobre la riqueza se denominan impuestos directos porque se aplican directamente a los individuos. Los impuestos sobre el consumo se denominan impuestos indirectos porque gravan a los individuos indirectamente al gravar sus transacciones.

La imposición en el mundo

Las Figuras 18-1 y 18-2 muestran la distribución de los ingresos fiscales entre estos diferentes tipos de impuestos en Estados Unidos y en otras naciones. El gobierno federal de Estados Unidos recibe una gran parte de sus ingresos fiscales (41,6 %) del impuesto sobre la renta individual y otra gran parte (34,8 %) de los impuestos sobre las nóminas. Una proporción más moderada procede del impuesto sobre la renta de las sociedades (15,1 %), con proporciones muy pequeñas procedentes de los impuestos sobre el consumo y otros impuestos (como los impuestos sobre la riqueza).

La distribución de las fuentes de ingresos fiscales en los niveles estatal y local es bastante diferente. Para los estados y las entidades locales, los ingresos procedentes del impuesto sobre la renta individual representan solo el 13,9 % de los ingresos totales. Las dos principales fuentes de ingresos para los gobiernos subnacionales de Estados Unidos son los impuestos sobre las ventas y los impuestos locales sobre la propiedad de viviendas y propiedades comerciales. Aunque algunos estados recaudan impuestos sobre las nóminas, los datos sobre los impuestos sobre las nóminas no están disponibles a nivel estatal. Para Estados Unidos en su conjunto, combinando todos los niveles de gobierno, la mayor cantidad de ingresos se recauda mediante la imposición sobre la renta, seguida de los impuestos sobre las nóminas y los impuestos sobre el consumo, mientras que los impuestos sobre la propiedad y los impuestos sobre la renta de las sociedades son aproximadamente iguales.

Las distribuciones son bastante diferentes en otras naciones. La Figura 18-2 muestra la distribución de los impuestos en dos países con sistemas fiscales muy diferentes, Noruega y Dinamarca, y después para el conjunto de naciones desarrolladas de la Organización para la Cooperación y el Desarrollo Económicos (OCDE) en su conjunto. En Noruega, los ingresos se recaudan casi por igual mediante impuestos sobre la renta, impuestos sobre las nóminas, impuestos sobre el consumo e impuestos sobre la renta de las sociedades, con una proporción muy pequeña de los ingresos procedente de la imposición sobre la propiedad. En Dinamarca, por el contrario, la mitad de los ingresos se recauda mediante impuestos sobre la renta individual, con una gran parte del resto procedente de los impuestos sobre el consumo; hay pequeñas proporciones procedentes de los impuestos sobre la renta de las sociedades, los impuestos sobre la propiedad y los impuestos sobre las nóminas. En promedio, las demás naciones desarrolladas de la OCDE tienen una proporción de ingresos fiscales procedentes de los impuestos sobre el consumo que es más de cuatro veces mayor que la de Estados Unidos y proporciones procedentes de los impuestos sobre las nóminas que son aproximadamente iguales.


Tipos impositivos medios y marginales

Dos conceptos clave describen el conjunto de tipos impositivos sobre la renta. El primero es el tipo impositivo marginal, el porcentaje del siguiente dólar de renta imponible que se paga en impuestos. Con un sistema como el de Estados Unidos, el tipo impositivo marginal aumenta con la renta; para quienes tienen una renta imponible inferior a 18.450 dólares, el tipo impositivo marginal es del 10 %, pero para quienes tienen una renta imponible superior a 464.850 dólares, el tipo impositivo marginal es del 35 %.

El segundo concepto es el tipo impositivo medio, el porcentaje de la renta total que se paga en impuestos, que se calcula como la razón entre los pagos totales de impuestos y la renta total.

El tipo impositivo medio de cualquier individuo es una media ponderada de los tipos marginales que el individuo paga a medida que avanza por la escala impositiva. Por ejemplo, supongamos que Josh tiene una renta bruta de 190.000 dólares y una renta imponible (después de ajustes, deducciones y exenciones) de 170.000 dólares (y ningún crédito fiscal). Su factura fiscal total es

(18.450 $ × 0,1) + (56.450 $ × 0,15) + (76.300 $ × 0,25) +
(38.800 $ × 0,28) = 34.651,50 $.

Calculamos la factura fiscal total de Josh haciéndolo avanzar por la escala de tipos marginales hasta llegar a su nivel de renta. El tipo marginal de Josh es del 28\% porque este es el tipo que paga sobre su siguiente dólar de renta. Su tipo impositivo medio es del 18 %, que es su factura fiscal (34.651,50 dólares) dividida por su renta bruta (190.000 dólares); esta es una media ponderada de todos los tipos marginales que Josh está pagando, donde las ponderaciones son la proporción de su renta en cada tramo impositivo (incluidos los 20.000 dólares de su renta que no están gravados y que, por lo tanto, se enfrentan a un tipo marginal cero).


\subsection{Diseño del Sistema Fiscal: Principios impositivos}
Los principios impositivos son en principio universales en gobiernos democráticos, recogidos generalmente en normas de rango legal elevado, sea la constitución o leyes mismas. Estos deben ser inspiradores de la actividad legislativa vinculada a los ingresos públicos y especialmente a los impuestos, además de orientar la actuación ejecutiva del proceso de recaudación ele impuestos. 

No debemos siempre ver las cosas desde una óptica nacional y por ello antes de nada observaremos los principios generales de la imposición desde una óptica teórica antes de mencionar los recogidos en nuestro ordenamiento jurídico.

De esta manera el debate parte de la coercitividad de la imposición. Los impuestos se recaudan de manera coactiva, en contra de los deseos de muchos individuos. Es por eso que los impuestos son una de las herramientas más poderosas que tienen los gobiernos a la hora de generar cambios de comportamiento en sus agentes. Es la herramienta de política económica más poderosa que existe (i.e: si se eximiese del pago de IVA a todo aquel que sallase a la comba. ¿qué ocurriría? ) y por eso se debe de analizar y diseñar con cui dado. Esta es por lo tanto la doble cara de los impuestos: por un lado son instrumento de política económica y por otro son elemento esencial para la financiación.

Los impuestos sirven al Gobierno para cumplir ciertos objetivos y es por el desarrollo de estos objetivos por el cual se establece una carga a los agentes económicos. Por lo tanto el diseño de un sistema tributario óptimo debe ir acompañado por unos objetivos económico básicos. No podemos analizar principios sin analizar estos objetivos. En una economía moderna, donde se persigue el bienestar de la sociedad, podemos decir que los objetivos económicos generales deben ser:
• Máxima libertad consistente con el bienestar de los demás.
• Niveles de vida óptimos de acuerdo a los recursos y tecnología disponibles.
• Crecimiento económico óptimo: sostenible, fuerte y estable.
• Distribución de la riqueza de acuerdo a los estándares de equidad dentro de la sociedad.

Es bajo estos objetivos generales de un Gobierno sobre los que se puede desarrollar un sistema impositivo acorde a los principios generales de la imposición. Aquí la palabra importante es sistema. El sistema tributario es un conjunto y no puede analizarse no de manera asilada ni de manera parcial. Es decir. siempre debe darse un enfoque de equilibrio general y de conjunto. Debemos tener en cuenta. en nuestra labor de economistas que no todos los países persiguen estos principios. Pensemos en economías como Guinea Ecuatorial o Mauritania, donde los conceptos de equidad o libertad no son perseguidos por el Gobierno. En dicho contexto no podremos usar la óptica de los principios generales para juzgar su sistema tributario. Los principios tributarios se han estudiado en economía desde el inicio de la disciplina. ADAM SMITH estableció en sus estudios lo que se denominan los principios canónicos de la tributación.\footnote{%
    Se recomienda entender estos principios canónicos en la óptica de la independencia de los EEUU y la revolución francesa.
}


• Equidad: La tributación debe ser proporcional al ingreso.
• Transparencia: El contribuyente debe saber cuánto paga.
• Conveniencia: Se debe suministrar alternativas de pago.
• Economía: Los impuestos deben ser acordes a los costes de recaudación. Si estos son excesivos no se deben establecer. 

Sin embargo. hoy en día los principios tributarios se configuran de una manera diferente de acuerdo a los objetivos anteriormente seí'ialados y a la configuración de Gobiernos modernos.

• Generalidad: Partiendo de la idea de que todos los individuos deben contribuir al sostenimiento del Gasto Público. se establece este principio. Asociado al mismo. la administración debe cumplir una serie de principios técnico tributarios, que facilitan el cumplimiento del ante1ior: o Sencillez Administrativa: Es decir se deben primar los impuestos de amplia base. D e manera que lo que grave el tributo se defina de forma genérica sin excepciones ni reglas complejas; así se reducen los costes de gestión y el fraude. o Continuidad y Flexibilidad: Las condiciones deben ser estables, sin cambios bruscos e inadvertidos, pero a su vez el tributo debe Sl'r capaz de captar las condiciones cambiantes de la economía. para que la recaudación se ajuste. Esto lleva a una preferencia por la progresividad y al grnvamen integro de rentas, asociado a bases amplias.

• Suficiencia. La recaudación de los impuestos debe ser suficiente como para atender los obj etivos del Gobierno representados en los estados de gasto de los Presupuestos Generales del Estado.

• Eficiencia. neutralidad y no neutralidad. El sistema tributario debe no provocar distorsiones sobre decisiones individuales de los agentes económicos. Pero a su vez un sistema tributario debe prestar incentivos o desincentivos que alteren el comportamiento de los agentes (i.e: impuestos ambientales o sobre tabaco).   Así pues, en primer lugar, un sistema tributario debe estar disefiado para favorecer la Eficiencia Económica, sin introducir distorsiones y ser Neutral sobre el comportamie1ito de los agentes. Esto se materializa en términos de Hacienda Pública minimizando el Exceso de Gravamen, es decir minimizando la perdida de bienestar que genera el establecimiento del impuesto.\footnote{Es un concepto asociado al efecto sustitución ele los impuestos.
} Por otro lado. la persecución de objetivos socialmente deseables no es patrimonio exclusivo del gasto. Los impuestos como figura de ingresos pueden estar diseñados para la obtención de dichos objetivos (protección de minusvalía. etc .... ) y por lo tanto estaremos en un caso de No neutralidad.

• Equidad. Al ser la imposición un fenómeno obligatorio, los impuestos deben ser percibidos como justos, la equidad será pues la medida de esa ·'justicia'' tributaria. De manera que la tributación se determine en virtud de criterios objetivos (renta o riqueza) y a su vez tenga en cuenta las circunstancias individuales. Así la cuestión clave del sistema tributario es tratar desigualmente lo desigual. La equidad es de dos tipos Vertical y Horizontal: Equidad Horizontal: Este concepto es ·'Los iguales deben de tributar igual''. Es decir, individuos que, acorde a la definición objetiva del hecho imponible, se encuentran en idéntica situación deben tributar igual. Este principio. aunque sencillo. queda modificado rápidamente y en oca<;iones no se cumplirá, pues dos individuos sujetos pasivos de un tributo, pueden tener situaciones personales diferentes ( hijos a cargo, por ejemplo) que hace que la tributación sea diferente pese a que aplicando el hecho imponible sean idénticos. Equidad vertical: Este c1iterio es complementario al anterior. De esta forma ·'los desiguales deben tributar de manera desigual ... Dicho de otra forma. aquellos que según el hecho imponible sean sujetos pasivos diforentes deberán pagar una cuota diferente. El principio económico detrás de este principio es el de marginalidad. La utilidad marginal de un euro, para una persona rica es menor que la utilidad de una persona pobre. De esta fomrn el impacto en la utilidad de los sujetos pasivos de la detracción de renta (representada esta en cestas de bienes. en términos microeconómicos) debe ser igual. Este principio se traslada al siguiente.

• Principio de beneficio y principio de capacidad de pago. Estos principios adaptan los conceptos de equidad de manera que son paralelos. Así el principio de beneficio es equivalente al de equidad vertical y propone un repa1to de la tributación de los desiguales y el principio de capacidad de pago que es una manifestación del principio de equidad horiZ0ntal, en tanto que busca establecer criterios para alcanzarla en torno a esa capacidad de pago. o El principio de beneficio indica que, dado que los impuestos se utilizan para ·que el gobierno proporcione servicios, la persona que más disfr��1ta de dichos beneficios debe contribuir más. Así, aquellos que reciben más servicios del gobierno (hospitales, colegios. carreteras. alumbrado) deben pagar más impuestos. Esto presenta un problema práctico. cómo estimar los beneficios obtenidos por los contribuyentes de los bienes y servicios públicos prestados. Ese argumento podría servir para imponer el peaje universal (quien quiera usar las caneteras que pague u n peaje quien no, no). Este principio en términos económicos es el más eficiente, pues evita el Free-Riding y muchos de los problemas que los bienes públicos presentan. Un ejemplo de tributo así definido son los impuestos Pigouvianos. Los problemas de estos tributos están en la medición de la utilidad de los bienes que sirve de indicador para cargar el tributo. ( Ver tema de Bienes Públicos y Externalidades).
o El principio de capacidad de pago, en cambio, señala que cacja individuo debe pagar según su capacidad, medida esta en vi1tud de un indicador de la misma. Supera todos los problemas del Principio de Beneficio, si bien incorpora los propios. De esta fonna los impuestos son redistributivos y alcanzan mayor equidad vertical. El problema surge de qué indicador usar para medir la capacidad de pago:

■ Riqueza. Presenta el problema de la productividad de la misma. Una persona con 5 pisos en propiedad tendrá mayor capacidad de pago que aquel que tiene un cuadro de Goya del mismo valor que los pisos.

■ Renta. El problema surge del esfuerzo realizado y de las circunstancias personales. Así una renta de 1000€ trabajando de sol a sol con 4 hijos, no es lo mismo que una renta de 1000€ obtenida por la lotería o el ·'sueldo Nescafé'' sin hijos.

■ Consumo. Dos trabajadores de idéntico sueldo. pueden gastar de manera diferente. De hecho, uno con muchas personas a su cargo gastará el 100% de su sueldo y uno con menos gastará el 80%. Si sólo nos fijamos en el consumo. el primero debería pagar, más que el segundo.

En el sistema jurídico español los principios tributarios se recogen tanto en la Constitución, como en la Ley General Tributaria. Pasamos por lo tanto a ver como se plasman estos principios generales en nuestro ordenamiento jurídico.

Así en el a11ículo 3.1 de la LGT se consagran los principios de:

• Justicia, supone que los impuestos deben ser equitativos. en los términos que hemos visto con anterioridad.

• Generalidad. supone que todos los individuos deben contribuir en tamo que son contribuyentes.

• Igualdad, inspirado en el artículo 1 4 de la Constitución. que supone que la carga tributaria debe repartirse equitativamente entre todas lc1s personas con igual capacidad económica.

• Progresividad y Equitativa distribución de la carga tributaria, suponen que la carga tributaria se reparta en función de la capacidad de pago, como hemos visto anteriormente. El sistema espaíiol. utilizará diferentes medidas de capacidad de pago dentro de cada tributo. Para el IV A: el consumo, para el I RPF, la renta etc ....

\begin{specialblock}{¿Cómo medir la progresividad de un impuesto?}
    La progresividad de los impuestos puede medirse de varias maneras
    
    Supongamos que un investigador ha conseguido calcular la carga real que cada persona soporta de un determinado impuesto —la incidencia económica, tal como se definió anteriormente—. El resultado final de este ejercicio suele consistir en caracterizar el impuesto como proporcional, progresivo o regresivo. La definición de un impuesto proporcional es sencilla: describe una situación en la que la relación entre los impuestos pagados y la renta es constante, independientemente del nivel de ingresos. Definir qué es un impuesto progresivo o regresivo no es tan sencillo y, desafortunadamente, las ambigüedades en la definición a veces generan confusión en el debate público. Una forma natural de definir estos términos es en función del tipo impositivo medio (average tax rate), es decir, la relación entre los impuestos pagados y la renta. Si el tipo impositivo medio aumenta con la renta, el sistema es progresivo; si disminuye, el impuesto es regresivo. La confusión surge porque algunas personas entienden la progresividad en términos del tipo impositivo marginal (marginal tax rate), es decir, la variación de los impuestos pagados con respecto a una variación de la renta. Para ilustrar la diferencia, considérese la siguiente estructura muy simple del impuesto sobre la renta. Cada individuo calcula su obligación tributaria restando 3.000 dólares a su renta y pagando una cantidad igual al 2\% de la diferencia. (Si la diferencia es negativa, el individuo recibe una subvención equivalente al 20 % de dicha cantidad). 
    
    La Tabla 14.1 muestra el importe del impuesto pagado, el tipo impositivo medio y el tipo impositivo marginal para varios niveles de renta. Los tipos medios aumentan con la renta. Sin embargo, el tipo impositivo marginal permanece constante en 0,2 porque, por cada dólar adicional ganado, el individuo paga 20 centavos adicionales, independientemente de su nivel de renta. Dos personas podrían discrepar acerca de la progresividad de este sistema tributario y, sin embargo, ambas tendrían razón de acuerdo con sus respectivas definiciones. Por ello, es muy importante dejar clara la definición utilizada cuando se emplean los términos regresivo y progresivo. A partir de este punto, asumiremos que estos términos se definen en función de los tipos impositivos medios. Medir el grado de progresividad de un sistema tributario es una tarea aún más difícil que definir la progresividad. Se han propuesto muchas alternativas razonables; aquí consideraremos dos medidas sencillas. La primera sostiene que cuanto mayor sea el incremento de los tipos impositivos medios a medida que aumenta la renta, más progresivo será el sistema. Algebraicamente, sean T_0 y T_1 las obligaciones tributarias reales (en contraposición a las legales) correspondientes a los niveles de renta I_0 e I_1, respectivamente, siendo I_1 > I_0. La medida de progresividad, v_1, es: v_1=\frac{\frac{T_1}{I_1}-\frac{T_0}{I_0}}{I_1-I_0} (14.1) Una vez que el analista calcula los valores de T_1 y T_0 y los sustituye en la ecuación (14.1), se considera más progresivo el sistema tributario que presente un valor más elevado de v_1. La segunda posibilidad consiste en afirmar que un sistema tributario es más progresivo que otro si la elasticidad de la recaudación tributaria con respecto a la renta (es decir, la variación porcentual de la recaudación dividida entre la variación porcentual de la renta) es mayor. En este caso, la expresión que debe evaluarse es v_2, definida como: v_2=\frac{\frac{T_1-T_0}{T_0}}{\frac{I_1-I_0}{I_0}} (14.2) Considérese ahora la siguiente propuesta: aumentar la obligación tributaria de cada persona en un 20 % del importe del impuesto que paga actualmente. Esta propuesta incrementaría la obligación tributaria de una persona que antes pagaba T_0 hasta 1,2 \times T_0, y la obligación que antes era T_1 hasta 1,2 \times T_1. El congresista A afirma que la propuesta hará que el sistema tributario sea más progresivo, mientras que el congresista B sostiene que no tendrá ningún efecto sobre la progresividad. ¿Quién tiene razón? Depende de la medida de progresividad utilizada. Si se sustituyen las expresiones 1,2 \times T_0 y 1,2 \times T_1 por T_0 y T_1, respectivamente, en la ecuación (14.1), el valor de v_1 aumenta un 20 %. Según esta medida, la propuesta incrementa la progresividad. En cambio, si se realiza la misma sustitución en la ecuación (14.2), el valor de v_2 permanece inalterado. (Tanto el numerador como el denominador se multiplican por 1,2, por lo que el efecto se cancela). La enseñanza es que incluso medidas de la progresividad que parecen muy intuitivas pueden conducir a conclusiones distintas. Una vez más, un debate público riguroso requiere que las definiciones utilizadas se hagan explícitas.
\end{specialblock}

• No confiscatoriedad, supone que no debe haber excesos por parte de los poderes públicos a la hora de detraer de las economías privadas los ingresos tributarios. Este debate es muy interesante pues se pude preguntar cuándo empieza la confiscatoriedad. Cuando se supera el tipo del 50% marginal o cuando se supera el 50% medio o en su caso sólo cuando se confisca toda la renta.














































































\subsection{Características: ¿Qué hace de un ingreso un \textit{impuesto}?}
Los impuestos los definimos como: \textit{Todo pago coactivo sin contraprestación directa que implica una transmisión coactiva de valores económicos del particulares al Estado en virtud de un mandato legal o constitucional}.

Esta definición recoge los tres elementos esenciales que componen un impuesto:
• Una obligación pecuniaria.
• Una obligación a favor de un ente público.
• Una obligación de carácter general o particular pero definida en vi1tud de un mandato legal.

Para hablar pues de un impuesto tenemos que tener una serie de elementos definidores o elementos esenciales, de acuerdo a la Constitución española, que no hace sino recoger los principios generales de Hacienda Pública que rigen a cualquier sistema y que se resumen en d siguiente esquema:

\textbf{INTRODUCIR ESQUEMA}

Pasamos a delinir cada elemento que nos ayudará posteriorm ente a estudiar y analizar cada tributo:\footnote{%
    Sirva este esquema como el discurso lógico para exponer esta parte. y como esquema de estudio general de los tributos concre10s que se analizan más ade.lanle en el temario.
}
• Suj eto activo: El Suj eto activo es aquel que ti<:;ne capacidad de fijar un impuesto, Depende de cada país. En el caso de España la habilitación es constitucional y viene recogida en el art. 133. que la repai1e de manera no simétrica entre Gobierno. Comunidades Autónomas y Entidades Locales. Siendo la principal responsabilidad del primero y del segundo, este siemrre en virtud del principio de autonomía financiera. La lógica de esta asimetría bebe de que la capllcidad de recolectar tributos es derivada de la soberanía nacional donde los individuos ceden al gobierno esta capacidad, este a su vez y en los casos cletenn inados por la Ley los puede ceder a Comunidades Autónomas y Entes Locales.

• Sujeto pasivo: Es todo sujeto económico con la obligación de contribuir por haber incurrido en el supuesto de hecho de la ley quc'imputa la obligación tributaria. De esta form a es un sujeto. físico o jurídico. nacional o extranjero. público o privado. Debemos diferenciar entre: o Contribuyente, al que la ley otorga la cargll tributarill. o Sustituto del mismo, como responsllble principal del mismo (por ej emplo. en las retenciones del trabajo el contribuyente es el trabaj ador, pero quien realiza el pago es el empresario). o Sujeto pasivo económico. es conceptualm ente aquel que efectiv¡¡meilte soporta el tributo, y puede ser alguno de los anteriores o un tercero.

• Hecho imponible: Parte de la definición de un objeto imponible. que entendemos como aquel elemento de la realidad sobre el que recae el tributo. Así pues, el hecho imponible es la materialización del objeto im ponible, supuestos los hechos jurídicos y económicos fij ados en la norma y que dan lugar a la ohligación trihutaria. Cabe la posibilidad de dos exclusiones: o No sujeción. No se da el hecho imponible, pese a darse el objeto imponible. porque no se cum ple algún supuesto económico jurídico fijado en la ley. o Exención. Se da el hecho imponible, pero la ley lo excluye específicamente (Por ejemplo. la exención de pago de I BI ll los bienes de la Iglesia Católica en España.)

• Base imponible: La base imponible es la cuantificación y val oración del hecho imponible. Generalmente es monetaria, pero también puede ser física en núm ero cigarrillos o KJ de energía. Se calcula por los siguientes métodos de estimación: o Directa; a partir de datos contables o la declaración vol untaria del contribuyente. o Indirecta. que se hace a partir de índices o módulos. o Objetiva, cuando no existen elementos de juicio, se realiza un análisis por testigos próximos.

• Base liquidable: es la base imponible, menos las reducciones a esta que marque la ley. Estas son minoraciones de la base imponible debido a motivos de política económica o simplemente diseño del impuesto. Es importante señalar que estamos ante ·'reducciones .. de la base. es decir que son motivos por los cuales el legislador considera que se debe reducir la carga fiscal por algún motivo asociado al hecho imponible. Ejemplo, El alquiler de imnueble a una persona para establecer su domicilio tiene una reducción para el arrendador del 60% de la renta generada en dicho ¡¡Jquiler.

• Tipo impositivo. La carga tributaria que recae sobre la base liquidable. puede ser de tipo proporcional o ad valorem. Asimismo. los tipos de gravamen se pueden juntar en una tarifa, que será un conjunto de tipos de gravamen que se aplicarán a un contribuyente. Así un contribuyente de JRPF del tramo superior paga todos los tramos de la tarifa un tipo impositivo diferente.

• Cuota integra: Es el equivalente de la cuota tributaria y es fruto de multiplicar la base liquidable por el tipo impositivo. Es en realidad lo que representa la carga o gravamen para el contribuyente en términos monetarios.

• Cuota líquida; es el equivalente de la deuda Tributaria, es decir lo que efectivamente se integra en las arcas del gobierno. pues se minora la Cuota Integra con una serie de Bonificaciones, Deducciones o se incrementa con Recargos. Estas deducciones son de especial importancia pues responden a medidas de política económica y tiene un importante efecto en el comportamiento de los agentes. Aquí debemos resaltar la diferencia de reducción en la base contra deducción en la cuota. que es el caso que estamos estudiando. Un ejemplo de deducción es la deducción por compra de vivienda habitual, a extinguir pero que supone un alivio fiscal de unos 3.000 millones de euros para los contribuyentes de IRPF.

• Cuota diferencial. Es la diferencia entre la cuota líquida y lo ya pagado mediante retenciones o ingresos a cuenta y que de facto supone el elemento de pago del impuesto.




\subsection{Clases de impuestos: ¿Qué impuestos podemos encontrar?}
Existen multitud de clasificaciones de impuestos, dependiendo de cada uno de los elementos definidores que hemos visto anteriormente. De manera reducida repasamos los más habituales. • Según Sujeto pasivo pueden ser Personales. que gravan a una persona fisica. o jurídica o reales, que gravan un bien. • Según como se defina la Base Imponible pueden ser analíticos. analizando cada fuente de base imponible o sintético��. sin valorar las fu entes. • Según la tarifa, y como esta opera en virtud de la base imponible, pueden ser proporcionales, la tarifa y la base crecen al mismo ritmo, progresivos. la tarifa crece más rápido que la base o regresivos, donde la tarifo crece menos que la base. • Según el tipo de gravamen que puede ser ad valorem o especifico. • Según el momento de cobro periódico o accidental.

Sin embargo, la clasificación más importante es aquella que gira en torno al hecho imponible y que vincularemos a la definición de contabilidad nacional hecha anteriormente.

• Directos. Son aquellos tributos que gravan una manifestación directa de capacidad de pago. Giran generalmente sobre la renta y la riqueza y por ello se asemejan, aunque no exclusivamente a los vistos anteriormente como ·'Impuestos corrientes sobre el ingreso, la riqueza etc.''. Pueden ser de carácter subjetivo, teniendo en cuenta las condiciones partic ulares del obligado tributario o de carácter objetivo.

• Indirectos. Son aquellos tributos que gravan una manifestación indirecta de capacidad de pago. Giran generalmente sobre flujos de riqueza. como puede ser la compra de inmuebles, o sobre usos de la renta y riqueza, como son el consumo. Suelen ser repercutibles y pueden implicar que la cuota se obtenga de un tercero diferente del obligado tributario.

• Cotizaciones sociales. Son contribuciones a un sistema de protección social y en cuanto a tales finalistas. Se asocian con el mercado de tra bajo (aunque nada impide que se carguen al mercado de capitales. aunque nunca se ha hecho) y la prestación de empleo por parte de un trabajador. El sujeto pasivo suele ser tanto el trabajador como el empleador y la determinación de las mismas puede ser directa. a raíz del sueldo. o indirecta a raíz de una base de cotización.








\clearpage
\section{Conclusión} 
\textbf{NOTA:} Esta concl usión sirva de ejemplo únicamente. El opositor puede preferir recapitular o utilizar una conclusión común parn este tema y los temas 4,5 y 6.


Los ingresos del Estado presentan una de las peculiaridades más interesantes cuando se valoran las funciones del m ismo. Por un l ado. los ingresos son instrumento de financiación de los gastos que se desan-ollan para realizar objetivos socialmente deseables. Por otro lado, los ingresos son a su vez una herramienta para modificar el comportamiento de los agentes económicos. igual que e l gasto.

Es esta doble naturaleza, financiera y finalista, la que hace que los ingresos y en particular los impuestos uno de los instrumentos más potentes de política económica que existen a disposición de los gobiernos. Esta herramienta a su vez es macroeconómica, en tanto que afecta a las fianzas del Estado y microeconómica, en tanto que afecta al comportamiento de los agentes individuales. Ninguna norma o actuación tiene el poder que tienen las tributarias.

Los ingresos públicos como componente macroeconómico, son generalmente tributarios. pero debemos dejar de lado los no tributarios que proviene de otras fuentes. en la mayoría de los casos rentas de propiedad. Estas otras fuentes no tributarias pueden a l canzar un volumen de tal nivel que permitan que toda actuación presupuestaria sea posible (ejemplo. los ingresos por petróleo).

Es por ello que a la hora de analizar una economía el conocimiento de las fuentes de ingresos y el tratamiento de las mismas es algo imprescindible no sólo para valorar elementos tan importantes como la posición fiscal, si no para entender el comportamiento de los agentes nacionales.

Los principios impositivos representan el debate en torno a la imposición. que siempre gira en torno al binomio: equidad - eficiencia, y señalan qué aspectos se deben tener en cuenta al diseñar las políticas de ingresos y las herramientas tributarias. Estos principios son esenciales. dado el poder de la política tributaria y d e ingreso, pues sin ellos nos e alcanzarían en ningún caso los objetivos socialmente deseados, pero no se debe olvidar que se debe analizar el conjunto del sistema tributario y utilizar los principios sobre el conjunto del mismo. Sin esta visión global es posible que se llegue a conclusiones equivocadas.

\subsection{Recapitulación}


\subsection{Extensiones y relación con otras partes del Temario}



\subsection{Opinión}

\subsubsection{Idea final (Salida o cierra)}

\clearpage
\pagenumbering{gobble}
\MyBibliography



\MyBibItem{DAWID y GATTI}{Chapter 2 - Agent-Based Macroeconomics}{2018}{https://doi.org/10.1016/bs.hescom.2018.02.006}


% ANEXOS
\clearpage










\end{document}